%% notes-tex/028-gene-essentiality/sections/5-families.tex
%% [[experiments.028-gene-essentiality]]
%% Moved here from the backbone section on 2026-09-21: every arm these subsections read is
%% defined in the module section, so they follow it.

\section{Which genes are missed, the two model families, and the genome-wide set\secstatus{todo}}
\label{sec:families}

The arms read here are the backbone cells of Section~\ref{sec:backbone} and the module
forms of Section~\ref{sec:module}, so every arm is defined before it is compared.

\subsection{Which genes the best cells misrank\secstatus{todo}}
\label{sec:hard-genes}

At 195 genes the AUROC is a sum over a few dozen misrankings, so the genes behind it
can be named. For every cell of nine arms, \texttt{hard\_genes.py} reads the
per-gene logits on the held-out genes and scores each gene's \emph{misrank}, the
fraction of the other class on the wrong side of it: for an essential gene the
nonessential genes scoring above it, for a nonessential gene the essential genes
scoring below it. One minus the class mean of the misrank is that class's share of the
AUROC, and averaging over an arm's cells separates the genes an arm gets wrong from the
ones a seed gets wrong (Figure~\ref{fig:hard-genes}, Table~\ref{tab:hard-genes}).

\tcfigfit[1.0]{figures/hard_genes_2026-09-21-10-23-58.pdf}{\textbf{Which held-out genes
  the best cells misrank.} Each arm has a letter, a color and a marker shape, shared by
  all three panels and listed in the key beside panel \textbf{c}; an arm's value for a
  gene is the mean over that arm's cells.
  \textbf{a}, The 31 essential held-out genes sorted by their mean misrank over arms,
  one marker per arm, the black line the mean over arms.
  \textbf{b}, The same for the 164 nonessential held-out genes.
  \textbf{c}, The Spearman correlation of the per-gene misrank between arms, with the
  value written in each cell.
  Arms: \textbf{A}, the pathway multi-hot with ProtT5 and the affine (round 22, three
  seeds); \textbf{B}, the round-12 fold models (five); \textbf{C}, the head over the
  solved flux on ProtT5 (rounds 21 and 24, six seeds); \textbf{D}, the same on the
  pathway input (rounds 23 and 25, six); \textbf{E}, that head over the cached GEM-box
  solve on the pathway input (round 27, six); \textbf{F}, the layer-free cell with the
  $\log(1 + |v|)$ vector in the affine (round 33, seeds 2700 to 2705); \textbf{G},
  ProtT5 with the affine (round 5, three); \textbf{H} and \textbf{I}, graph convolution
  on GO at 400 epochs, three seeds and six. Written by
  \texttt{experiments/028-gene-essentiality/scripts/hard\_genes.py}.}{fig:hard-genes}

%% notes-tex/028-gene-essentiality/tables/hard_genes.tex
%% GENERATED by experiments/028-gene-essentiality/scripts/notes_tex_tables.py
%% SOURCE: results/hard_genes.json
\begin{table}[H]
  \centering
  \scriptsize
  \caption[The hardest held-out genes]{The fifteen held-out genes of each class the best cells misrank most, sorted by the mean misrank over the arms. A gene's misrank in one cell is the fraction of the other class on the wrong side of it (for an essential gene, the nonessential genes scoring above it; for a nonessential gene, the essential genes scoring below it), averaged over the arm's cells; one minus the class mean of the misrank is that class's share of the AUROC. FBA is the knockout-to-wild-type growth ratio on YPD; pathways and GO the gene's counts in the pathway multi-hot and the propagated GO vector. Arms: A1: pathway + ProtT5 + affine (round 22); A2: indicator + ProtT5 + affine, folds (round 12); A3: flux head, ProtT5 (rounds 21, 24); A4: flux head, pathway + ProtT5 (rounds 23, 25); A5: ProtT5 + affine (round 5). Written by \texttt{hard\_genes.py} from the per-gene prediction dumps.}
  \label{tab:hard-genes}
  \begin{tabular}{@{}lrrrrrrrr@{}}
    \hline
    gene & A1 & A2 & A3 & A4 & A5 & FBA & pathways & GO \\
    \hline
    \multicolumn{9}{@{}l}{\emph{essential genes}} \\
    YDR373W & 0.44 & 0.28 & 0.51 & 0.47 & 0.59 & 1.00 & 0 & 25 \\
    YKL004W & 0.48 & 0.25 & 0.50 & 0.22 & 0.66 & 1.00 & 1 & 39 \\
    YML125C & 0.37 & 0.46 & 0.40 & 0.41 & 0.41 & 1.00 & 0 & 27 \\
    YAL038W & 0.13 & 0.42 & 0.62 & 0.29 & 0.49 & 1.00 & 2 & 81 \\
    YBL030C & 0.30 & 0.22 & 0.46 & 0.33 & 0.57 & 1.00 & 0 & 68 \\
    YMR296C & 0.14 & 0.31 & 0.36 & 0.42 & 0.35 & 1.00 & 1 & 36 \\
    YIL048W & 0.17 & 0.41 & 0.19 & 0.40 & 0.35 & 1.00 & 0 & 52 \\
    YDR062W & 0.15 & 0.15 & 0.40 & 0.43 & 0.27 & 1.00 & 1 & 36 \\
    YDR050C & 0.11 & 0.14 & 0.19 & 0.07 & 0.40 & 1.00 & 2 & 85 \\
    YBR011C & 0.10 & 0.10 & 0.16 & 0.33 & 0.14 & 1.00 & 0 & 18 \\
    YDL120W & 0.12 & 0.10 & 0.07 & 0.14 & 0.21 & 1.00 & 0 & 50 \\
    YMR281W & 0.07 & 0.23 & 0.06 & 0.17 & 0.12 & 1.00 & 0 & 64 \\
    YDR454C & 0.27 & 0.15 & 0.03 & 0.08 & 0.08 & 0.00 & 5 & 65 \\
    YGR277C & 0.05 & 0.04 & 0.03 & 0.33 & 0.05 & 1.00 & 1 & 64 \\
    YDL015C & 0.07 & 0.09 & 0.05 & 0.22 & 0.06 & 1.00 & 1 & 31 \\
    \multicolumn{9}{@{}l}{\emph{nonessential genes}} \\
    YGL012W & 0.26 & 0.83 & 0.79 & 0.72 & 0.72 & 0.00 & 2 & 48 \\
    YJR121W & 0.73 & 0.67 & 0.35 & 0.44 & 0.68 & 0.34 & 0 & 121 \\
    YOL049W & 0.20 & 0.34 & 0.74 & 0.65 & 0.77 & 0.00 & 2 & 27 \\
    YHR106W & 0.61 & 0.41 & 0.57 & 0.57 & 0.45 & 1.00 & 1 & 30 \\
    YBR121C & 0.65 & 0.52 & 0.32 & 0.42 & 0.62 & 1.00 & 0 & 66 \\
    YLR450W & 0.67 & 0.67 & 0.32 & 0.49 & 0.37 & 1.00 & 3 & 63 \\
    YGR106C & 0.54 & 0.48 & 0.59 & 0.35 & 0.52 & 1.00 & 0 & 30 \\
    YLR377C & 0.65 & 0.43 & 0.31 & 0.59 & 0.39 & 1.00 & 3 & 26 \\
    YDR315C & 0.45 & 0.41 & 0.58 & 0.32 & 0.54 & 1.00 & 0 & 73 \\
    YGR192C & 0.63 & 0.67 & 0.24 & 0.54 & 0.16 & 1.00 & 3 & 125 \\
    YJR009C & 0.65 & 0.52 & 0.25 & 0.54 & 0.17 & 1.00 & 3 & 113 \\
    YJL052W & 0.62 & 0.47 & 0.24 & 0.55 & 0.16 & 1.00 & 3 & 107 \\
    YAL054C & 0.65 & 0.24 & 0.30 & 0.55 & 0.31 & 1.00 & 2 & 76 \\
    YLR143W & 0.47 & 0.23 & 0.36 & 0.43 & 0.52 & 1.00 & 0 & 16 \\
    YBR111C & 0.37 & 0.35 & 0.35 & 0.37 & 0.52 & 1.00 & 0 & 23 \\
    \hline
  \end{tabular}
\end{table}


The solve's own call splits the essential genes. Nineteen of the 31 essential held-out
genes have a knockout growth ratio of 1.0 on YPD, the solve calling them dispensable,
and thirteen of the fifteen hardest are among those nineteen, while nine of the
fifteen easiest have a ratio of 0. The hardest essential genes are the ones FBA misses,
and every model misses them with it: AUR1, FRQ1, PET9 (the ADP/ATP carrier, covered in
the GEM by AAC1 and AAC3), NEO1, PGA3, CDC19 (pyruvate kinase, covered by its paralog
PYK2), LCB1 and LCB2 (serine palmitoyltransferase), IPP1, TPI1. The hardest
nonessential genes are metabolic enzymes the models call essential: ERG4, HMG2, GSH2,
the three glyceraldehyde-3-phosphate dehydrogenase isozymes TDH3, TDH2 and TDH1, ATP2,
the two phosphofructokinase subunits PFK1 and PFK2, TRR2, ERG2 and FBP1, twelve of the
fifteen with a ratio of 1.0 and most with 26 to 125 GO terms. Isozyme cover explains a
minority of the missed essential genes. Of the 19 essential held-out genes the solve
calls viable, the gene-protein-reaction rules (\texttt{isozyme\_affine\_store.py},
rules in disjunctive normal form) give an alternative gene set for every reaction of
only four, PET9, CDC19, PGA3 and YDR302W; the other fifteen, AUR1, NEO1, LCB1 and
LCB2 among them, each have at least one reaction that only they catalyze, which their
deletion closes and which the growth program does not need. So the larger failure is a
reaction the biomass objective does not require, and the smaller one a paralog that
covers the gene in the solve; among the nonessential genes it is the reverse case, a
paralog of an essential metabolic function that the annotation vectors make look
essential. The arms agree on the genes they fail: the two
affine cells at a Spearman of 0.81, the flux head on ProtT5 with the affine on ProtT5
at 0.88, the two flux heads on the pathway input with each other at 0.77, and those
two least aligned with the rest (0.43 to 0.76), the arms whose means carry an unfit
seed; the layer-free cell of round 33 sits with the transformer family, at 0.82 with
the pathway affine and 0.81 with the cached-solve head, and at 0.28 with the GO
convolution.
Hypothesis, untested: a paralog-aware input, the GPR disjunction with an expression
weight or the isozyme count of a gene's reactions, is what would move the fourteen; the
solve's own vector cannot, since the solve reroutes through the paralog.

Graph convolution on GO fails on different genes. Its per-gene misrank correlates
with the transformer arms' at 0.15 to 0.42 where those arms correlate with one another
at 0.4 to 0.9, the same on three seeds and on six, and its class mean misrank on the
essential genes is 0.091 against the
pathway cell's 0.110. On the hardest essential genes it halves the misrank the solve
leaves: FRQ1 from 0.44 to 0.17, AUR1 from 0.48 to 0.22, PGA3 from 0.37 to 0.11, LCB1
from 0.14 to 0.03, LCB2 from 0.15 to 0.01, while it does worse on NEO1 (0.17 to 0.60)
and PET9 (0.30 to 0.38). Among the nonessential genes it repairs twelve the pathway
cell misranks by more than 0.3 (TDH1, ACS1, ATP2, FBP1, TRR2 among them) and breaks six
the pathway cell had right. Two models that read the same genes wrong for different
reasons are the case for averaging them.

\subsection{The essential genes every predictor ranks low\secstatus{todo}}
\label{sec:consensus-missed}

If every predictor failed on the same essential genes, the shared failures would mark a
branch of function that no input describes. The released split holds out 31 essential
genes, too few to test that, so the intersection is taken on the ten repeated outer
splits of Section~\ref{sec:round75}, which hold out 874 of the 1,121 labeled genes at
least once, 140 of them essential. Five predictors are read there: the transformer with
the single-experiment module, the same transformer with no metabolic column, graph
convolution on GO, the DCell hierarchy with a shared input
(Section~\ref{sec:round108}), and FBA alone as one minus the YPD growth ratio. A gene's
misrank under one predictor is averaged over its seeds and over the splits that hold it
out, and its consensus misrank is the lowest of the four learned predictors.

The intersection is nearly empty. The mean misrank of the essential genes is 0.133 for
the module transformer, 0.193 without metabolism, 0.132 for graph convolution, 0.190 for
the DCell hierarchy and 0.290 for FBA, and each learned predictor ranks 7 to 15 essential
genes below the median nonessential gene of their block. No essential gene is below that
median under all four: 111 of the 140 are below it under none, 16 under one, 9 under two
and 4 under three. The consensus misrank averages 0.055, under half of the best single
predictor's mean, and only one gene reaches 0.5, the biotin transporter \textit{VHT1}
(Table~\ref{tab:consensus-missed}). The predictors disagree on which genes are hard: the
per-gene misrank of the essential genes correlates at a Spearman of 0.78 between the two
transformers, 0.40 between the module transformer and graph convolution, 0.43 between
graph convolution and the DCell hierarchy although both read the same GO multi-hot, and
0.08 between graph convolution and FBA.

What the residual shares is a function, and the depth of a gene's annotation does not
explain it. Across the 140 essential genes the consensus misrank has a Spearman
correlation of $-0.002$ with the count of propagated GO biological-process terms and
$-0.022$ with physical degree, and a term count under the class median is not
enriched. Of 179 terms and features tested by the Mann-Whitney rule of
Section~\ref{sec:hard-genes}, 11 pass $q < 0.1$. The largest group is the 70 essential
genes whose FBA growth ratio is 1.0, at a consensus misrank of 0.082 against 0.028 for
the rest (effect 0.78). Among GO terms the strongest are organic substance transport (10
genes, 0.190 against 0.045, effect 0.83), cellular component organization (14 genes,
0.120 against 0.048) and monoatomic ion transport (5 genes, 0.184 against 0.050). Test
blocks overlap across splits, so these $q$-values rank terms and are not calibrated error
rates. The top of the table is transporters and cofactor supply: \textit{VHT1},
\textit{HIP1}, \textit{RIM2}, \textit{THI80}, \textit{NCE103}. One reading is that these are essential only given the strain the deletion collection
was made in, BY4741, which carries \textit{his3}, \textit{leu2}, \textit{met15} and
\textit{ura3} deletions and lives on imported histidine, leucine, methionine and
uracil, while every FBA column here deletes one gene from the wild-type model. A probe
tested that (\texttt{fba\_background\_probe.py}): delete the four background genes,
take that strain as the reference, and recompute the knockout growth ratio of the 1,121
labeled genes. The background does not explain the residual at the resolution of the
model. On YPD 63 ratios move by more than 0.01 and none of them is an essential gene
moving toward lethal; the one gene that becomes lethal is the uracil permease
\textit{FUR4}, labeled nonessential, and the count of essential genes at a ratio of 1.0
goes from 77 to 83. FBA's AUROC over the labeled genes moves from 0.717 to 0.713 on YPD
and from 0.718 to 0.714 on SC, with ratios rounded to six decimals because solver noise
near a ratio of 1.0 reorders ties between identical runs. \textit{HIP1} and \textit{VHT1} stay at 1.0 because the
model's rules let other permeases carry their substrates. Hypothesis, untested: those
alternatives are repressed or too weak in rich medium, which is regulation and
transporter affinity that neither the model nor GO encodes. The probe also found a
defect in the modeled medium: YPD as defined here opens no uracil uptake, so a
\textit{ura3} deletion alone takes the model's YPD growth from 0.535 to zero (SC, which lists uracil, is unaffected), while
BY4741 grows on YPD. The probe opens uracil uptake in both of its arms; the stored
solves of the module were computed on the wild-type model, where the missing uptake
does not bind. \textit{ILV3} and \textit{ILV5}, which FBA ranks
correctly and the learned predictors do not, are the converse case of a label the
network explains and the annotation vectors do not.

%% notes-tex/028-gene-essentiality/tables/consensus_missed.tex
%% GENERATED by experiments/028-gene-essentiality/scripts/notes_tex_tables.py
%% SOURCE: results/consensus_missed.json
\begin{table}[H]
  \centering
  \footnotesize
  \caption[Essential genes every predictor ranks low]{The twelve essential genes with the highest consensus misrank over the ten repeated outer splits. A misrank is the share of the nonessential genes of the same test block that score above the gene, averaged over seeds and over the splits that hold the gene out (held), so 0 is a perfect rank and 0.5 the median nonessential gene. The consensus is the lowest of the four learned predictors, in bold: the transformer with the single-experiment module (nine seeds per split), the same transformer with no metabolic column (nine), graph convolution on GO (six) and the DCell hierarchy with a shared input (three). FBA scores one minus the YPD growth ratio and is not part of the consensus. Ratio is the FBA knockout growth ratio on YPD, GO the count of propagated biological-process terms, degree the physical interaction degree. Written by \texttt{notes\_tex\_tables.py} from \texttt{consensus\_missed.py}.}
  \label{tab:consensus-missed}
  \begin{tabular}{@{}llrrrrrrrrr@{}}
    \hline
    gene & systematic & held & module & no metab. & GO conv. & DCell & FBA & ratio & GO & degree \\
    \hline
    \textit{VHT1} & YGR065C & 2 & \textbf{0.50} & 0.58 & 0.75 & 0.68 & 0.57 & 1.00 & 14 & 13 \\
    \textit{NCE103} & YNL036W & 2 & 0.78 & 0.84 & 0.78 & \textbf{0.45} & 0.56 & 1.00 & 18 & 12 \\
    \textit{ACS2} & YLR153C & 3 & 0.38 & 0.38 & \textbf{0.32} & 0.51 & 0.55 & 1.00 & 60 & 80 \\
    \textit{HIP1} & YGR191W & 1 & 0.82 & 0.83 & 0.66 & \textbf{0.32} & 0.56 & 1.00 & 17 & 17 \\
    \textit{SDH3} & YKL141W & 4 & 0.40 & 0.99 & 0.59 & \textbf{0.31} & 0.06 & 0.94 & 38 & 19 \\
    \textit{ILV5} & YLR355C & 3 & 0.41 & 0.64 & \textbf{0.29} & 0.82 & 0.06 & 0.96 & 27 & 76 \\
    \textit{ARH1} & YDR376W & 3 & 0.70 & 0.51 & \textbf{0.28} & 0.57 & 0.56 & 1.00 & 26 & 10 \\
    \textit{RIM2} & YBR192W & 2 & 0.32 & 0.58 & 0.57 & \textbf{0.27} & 0.56 & 1.00 & 20 & 17 \\
    \textit{THI80} & YOR143C & 1 & 0.30 & 0.47 & 0.29 & \textbf{0.27} & 0.56 & 1.00 & 37 & 13 \\
    \textit{ILV3} & YJR016C & 6 & \textbf{0.24} & 0.42 & 0.30 & 0.32 & 0.06 & 0.96 & 22 & 35 \\
    \textit{CYR1} & YJL005W & 2 & 0.28 & \textbf{0.21} & 0.28 & 0.24 & 0.56 & 1.00 & 12 & 54 \\
    \textit{PET9} & YBL030C & 1 & \textbf{0.21} & 0.26 & 0.36 & 0.73 & 0.56 & 1.00 & 36 & 146 \\
    \hline
  \end{tabular}
\end{table}


\subsection{Averaging the two families\secstatus{todo}}
\label{sec:family-ensembles}

The average is a post hoc read of predictions already made, with nothing fitted: each
family's per-cell logits on the held-out genes are standardized and averaged over its
seeds, the families are averaged with equal weight, and the mean is scored once
(Table~\ref{tab:family-ensembles}). The GO convolution and the pathway cell with the
affine average to 0.933 [0.89, 0.97] from 0.909 and 0.891 alone; the GO convolution
and the flux head on the pathway input, six seeds with two unfit among them, average
to 0.941 [0.90, 0.98] from 0.909 and 0.911; all three to 0.935. With the six-seed
families of the later rounds, the GO convolution on six seeds (0.905 as a family) and
the scalable module form of round 27 (0.918 as a family, the highest single family)
average to 0.946 [0.91, 0.98], and with the pathway affine beside them to 0.939; the
layer-free form of round 33 (0.910 as a family) with the six-seed GO convolution
averages to 0.948 [0.91, 0.98], and the same cell over the interior-point store
(round 36, 0.912 as a family) to 0.944. Every average tried is in the table, so 0.948
is the maximum of fifteen post hoc choices and carries that bias, and the held-out
genes were used to choose nothing else. It says
what the per-gene
analysis said: the two families carry different information about the same genes, and
the sum of the annotation read through the graph and the solve read through the
transformer is more than either. Under the reported protocol the same average holds
(Table~\ref{tab:family-ensembles-folds}): the five GO fold models and the five
pathway flux-head fold models, each family 0.917 and 0.897 as its own ensemble,
average to 0.949 [0.91, 0.98]; the GO fold models with the reported cell's fold models
(0.893 alone) to 0.944; all three to 0.945. With the scalable form's five fold models
(round 27; 0.912 as its own ensemble and 0.922 as a standardized family, the highest
single family under this protocol) the GO fold models average to 0.954 [0.92, 0.98],
and with the reported cell's fold models beside them to 0.947; the layer-free form's
five fold models (round 34; 0.920 as their own ensemble and as a family) with the GO
fold models average to 0.953 [0.92, 0.98]. Eleven averages tried, all listed, the
best the maximum of eleven. The interval of 0.954 excludes the reported
0.893 point, and the two families' own intervals overlap it, which is what two models
that fail on different genes produce when averaged and what neither produces alone.

%% notes-tex/028-gene-essentiality/tables/family_ensembles.tex
%% GENERATED by experiments/028-gene-essentiality/scripts/notes_tex_tables.py
%% SOURCE: results/family_ensembles.json
\begin{table}[H]
  \centering
  \footnotesize
  \caption[Averages of the two model families]{Post hoc averages of existing predictions on the 195 held-out genes: each family's per-cell logits are standardized and averaged over its seeds, families are averaged with equal weight, and the mean is scored once with a 2,000-draw gene bootstrap. Nothing is fitted and every average tried is listed, so the best row is the maximum of seven post hoc choices, not a selected model. Written by \texttt{family\_ensembles.py}. Bold marks the highest value in a column.}
  \label{tab:family-ensembles}
  \begin{tabular}{@{}p{118mm}rl@{}}
    \hline
    average of & cells & AUROC [95\% CI] \\
    \hline
    GO convolution (400 epochs, 3 seeds) & 3 & 0.909 [0.85, 0.95] \\
    pathway + ProtT5 + affine (round 22, 3 seeds) & 3 & 0.891 [0.83, 0.94] \\
    flux head, pathway + ProtT5 (rounds 23 and 25, 6 seeds) & 6 & 0.911 [0.86, 0.95] \\
    GO convolution (400 epochs, 3 seeds) $+$ pathway + ProtT5 + affine (round 22, 3 seeds) & 6 & 0.933 [0.89, 0.97] \\
    GO convolution (400 epochs, 3 seeds) $+$ flux head, pathway + ProtT5 (rounds 23 and 25, 6 seeds) & 9 & \textbf{0.941} [0.90, 0.97] \\
    GO convolution (400 epochs, 3 seeds) $+$ flux head, ProtT5 (rounds 21 and 24, 6 seeds) & 9 & 0.923 [0.87, 0.97] \\
    GO convolution (400 epochs, 3 seeds) $+$ pathway + ProtT5 + affine (round 22, 3 seeds) $+$ flux head, pathway + ProtT5 (rounds 23 and 25, 6 seeds) & 12 & 0.935 [0.89, 0.97] \\
    \hline
  \end{tabular}
\end{table}


%% notes-tex/028-gene-essentiality/tables/family_ensembles_folds.tex
%% GENERATED by experiments/028-gene-essentiality/scripts/notes_tex_tables.py
%% SOURCE: results/family_ensembles_folds.json
\begin{table}[H]
  \centering
  \footnotesize
  \caption[Averages of the two model families, fold protocol]{The same post hoc averages under the reported protocol: each family is its five fold models (one seed per fold, the cross-validation files of Section~\ref{sec:backbone}), standardized per fold and averaged, families averaged with equal weight, scored once on the 195 held-out genes with a 2,000-draw gene bootstrap. Six averages tried, all listed. Written by \texttt{family\_ensembles.py --folds}. Bold marks the highest value in a column.}
  \label{tab:family-ensembles-folds}
  \begin{tabular}{@{}p{118mm}rl@{}}
    \hline
    average of & cells & AUROC [95\% CI] \\
    \hline
    GO convolution, five folds (400 epochs) & 5 & 0.917 [0.87, 0.95] \\
    indicator + ProtT5 + affine, five folds (round 12) & 5 & 0.897 [0.84, 0.94] \\
    flux head, pathway + ProtT5, five folds (round 25) & 5 & 0.897 [0.84, 0.94] \\
    GO convolution, five folds (400 epochs) $+$ indicator + ProtT5 + affine, five folds (round 12) & 10 & 0.944 [0.90, 0.98] \\
    GO convolution, five folds (400 epochs) $+$ flux head, pathway + ProtT5, five folds (round 25) & 10 & \textbf{0.949} [0.91, 0.98] \\
    GO convolution, five folds (400 epochs) $+$ indicator + ProtT5 + affine, five folds (round 12) $+$ flux head, pathway + ProtT5, five folds (round 25) & 15 & 0.945 [0.91, 0.97] \\
    \hline
  \end{tabular}
\end{table}


One model over both inputs does not recover the average. The deletion's solved flux
as the round-33 cell encodes it, $\log(1 + |v|)$ over 4,131 reactions, put on the
node beside the GO multi-hot and convolved over the union graph at 400 epochs reads
0.900, 0.912 and 0.914 (mean 0.909), the GO-alone rows' 0.909, and 0.910 with the
ratios beside them; alone through the graph it reads 0.773 and on the MLP 0.709, and
GO with it on the MLP 0.851 (Table~\ref{tab:graph-go-solve}). The convolution takes
nothing from the solve columns that the average of the two families took (0.948),
so what the average adds is not a matter of the graph model seeing the solve; the
solve's ranking, as the transformer's affine reads it, is not what a convolution
over these graphs extracts from the same columns.

The ten repeated splits agree (job 2684, three seeds per split). Against graph
convolution on GO alone, the same convolution with the solve columns on the node reads
$-0.009$ (higher on one of ten splits), with the four ratios beside them $-0.010$ (one
of ten), the MLP on GO and the solve columns $-0.013$ (four of ten), and the solve
columns alone through the graph $-0.189$ (none of ten). As node features the module does not carry to the graph family.

As an affine on the logit it does (job 2691, three seeds per split). The columns stay out
of the message passing and a zero-initialized linear map over them is added to each
node's logit, the form the transformer uses. Against the same convolution on GO alone,
rerun in the same job, the YPD ratio alone reads $+0.011$ over the ten splits (higher on
ten of ten), the single-experiment recipe of ratio, vector and counts $+0.018$ (nine of
ten) and the four-media recipe $+0.021$ (nine of ten); on the released split the three
read 0.914, 0.911 and 0.910 against 0.908. At six seeds per split (job 2693) the three
read $+0.009$ (ten of ten), $+0.017$ (seven of ten) and $+0.019$ (seven of ten), and
0.914, 0.911 and 0.911 against 0.905 on the released split. The interface that carries the module across
model families is therefore the affine on the logit, with the solve computed offline, and
it needs nothing from the architecture it is added to.

Giving both families the module does not improve their average. With the
single-experiment transformer of Section~\ref{sec:round75} and the graph family carrying
the same recipe as an affine, the fixed-in-advance average reads 0.897 over the ten
splits against 0.907 with the plain GO convolution as the partner, although the graph
family alone rises from 0.871 to 0.887. The average gains from the two families failing
on different genes, and handing them the same columns makes them fail more alike. The
fold protocol on the released genes says the same (job 2695, three seeds per fold): the
graph family with the affine reads 0.930 [0.88, 0.97] as a five-fold ensemble against
0.917 without, and averaged with the transformer's single-experiment fold models it
reads 0.941 against 0.953 for the plain graph family with the same partner.

The worst case for a module carried to the whole database is a record whose deletion
touches no reaction, which reads the constant wild-type row. The genome-wide held-out set
of Section~\ref{sec:round30} is entirely that case: none of its 4,674 genes is in the
solve store and every one has a growth ratio of 1.0, so the affine adds the same constant
to every held-out logit and cannot reorder them by itself. It still does not hurt (job
2696, three seeds, the same training and validation genes as the released split): graph
convolution on GO reads 0.702 alone and 0.767 with the single-experiment recipe on its
logit, higher on three of three seeds, and 0.704 with the YPD ratio alone. Since the
held-out logits differ only by the convolution's weights, the gain is in what the
convolution learned. Hypothesis, untested: the affine explains the metabolic genes'
labels during training, so the convolution's weights and its validation-selected epoch
(87 to 93 against 382 to 397 without the affine) fit the remaining genes rather than
the metabolic ones. The graph family's best genome-wide input, ProtT5 on the node
(Section~\ref{sec:round30}), reads the same way (job 2697, six seeds): on the released
split the single-experiment affine takes it from 0.859 to 0.891, higher on six of six
seeds, and on the 4,674 genome-wide genes from 0.751 to 0.766, six of six, with the YPD
ratio alone worth $+0.014$ and $+0.004$.

%% notes-tex/028-gene-essentiality/tables/graph_baselines_go_solve.tex
%% GENERATED by experiments/028-gene-essentiality/scripts/notes_tex_tables.py
%% SOURCE: results/graph_baselines_go_solve.json
\begin{table}[H]
  \centering
  \scriptsize
  \caption[One model over the annotation and the solve]{The deletion's pFBA flux vector as $\log(1 + |v|)$ (the round-33 encoding, 4,131 columns, zero for a gene the store does not cover) as a node input: alone, beside the GO multi-hot, and beside GO with the four FBA ratios, on the MLP (200 epochs) and on graph convolution over the union of the nine graphs (400 epochs, as Table~\ref{tab:graph-go400}). The protocol of Table~\ref{tab:graph-baselines} otherwise; seeds 2700 to 2702. Bold marks the highest value in a column.}
  \label{tab:graph-go-solve}
  \begin{tabular}{@{}llp{30mm}rrrrll@{}}
    \hline
    model & graph & inputs & dim & params & mean & sd & seeds & selected epochs \\
    \hline
    MLP & none & solve log1p & 4,131 & 268,673 & 0.709 & 0.012 & 0.716, 0.692, 0.717 & 86, 17, 93 \\
    GCN & union of nine & solve log1p (400 epochs) & 4,131 & 268,673 & 0.773 & 0.002 & 0.771, 0.772, 0.775 & 393, 393, 388 \\
    MLP & none & GO $+$ solve log1p & 6,493 & 419,841 & 0.851 & 0.000 & 0.851, 0.852, 0.851 & 23, 25, 26 \\
    GCN & union of nine & GO $+$ solve log1p (400 epochs) & 6,493 & 419,841 & 0.909 & 0.006 & 0.900, 0.912, 0.914 & 376, 397, 397 \\
    MLP & none & GO $+$ FBA ratios $+$ solve log1p & 6,497 & 420,097 & 0.851 & 0.001 & 0.851, 0.852, 0.850 & 25, 27, 26 \\
    GCN & union of nine & GO $+$ FBA ratios $+$ solve log1p (400 epochs) & 6,497 & 420,097 & \textbf{0.910} & 0.001 & 0.911, 0.912, 0.909 & 396, 397, 397 \\
    \hline
  \end{tabular}
\end{table}


The transformer side of the same question is round 47 (jobs 2434 and 2438;
Table~\ref{tab:w63}, Figure~\ref{fig:w63}): the GO multi-hot on the node beside the
pathway block and ProtT5 with the $\log(1 + |v|)$ vector in the affine, six seeds, and
the same with GO in place of the pathway block, beside a same-pack control with GO on
the node and the four ratios only. GO costs the transformer what it cost in round 26,
with the solve in the affine as with the ratios: 0.866 (sd 0.009; pooled 0.872 [0.81,
0.93]) against the round-33 cell's 0.908 and 0.912, with GO in place of the pathway
block 0.846 (pooled 0.855), and the control 0.852 (pooled 0.854), the round-26 cells'
0.849; the GO cells select at epochs 42 to 194 where the round-33 cell selects at 33
to 58. So neither family takes the other's input as the input enters it now, the
convolution nothing from the solve columns on the node and the transformer a cost
from GO on the node; whether GO enters the transformer the way the solve does, on the
logit through the affine rather than on the node, is round 50's read.

%% notes-tex/028-gene-essentiality/tables/round_w63.tex
%% GENERATED by experiments/028-gene-essentiality/scripts/notes_tex_tables.py
%% SOURCE: results/w63_readout.json
\begin{table}[H]
  \centering
  \scriptsize
  \caption[Round 47, the GO multi-hot beside the layer-free solve in one transformer]{Round 47, the GO multi-hot beside the layer-free solve in one transformer. The round-33 cell and the round-26 GO cells (GO beside ProtT5, and GO beside the pathway block and ProtT5, the four ratios only) as the comparators, against GO beside the pathway block and ProtT5 with the $\log(1 + |v|)$ vector in the affine, GO in place of the pathway block with the vector, and GO beside the pathway block and ProtT5 without the vector run in the same pack. Batch 32, class weight, 200 epochs; seeds 2700 to 2705 for the round-33 cell, 2700 to 2702 for the rest. Bold marks the highest value in a column. $\uparrow$ ($\downarrow$) marks an arm whose held-out AUROC is above (below) the first arm's on every one of the same three seeds.}
  \label{tab:w63}
  \begin{tabular}{@{}p{52mm}rrlllr@{}}
    \hline
    arm & mean & sd & seeds & pooled [95\% CI] & selected epochs & val step \\
    \hline
    the vector as log(1 $+$ |v|), default initialization & \textbf{0.908} & 0.008 & 0.896 to 0.919 (6) & \textbf{0.912} [0.87, 0.95] & 33 to 58 & 0.0043 \\
    GO multi-hot (propagated, 2,362) $+$ ProtT5 backbone $+$ FBA ratios & 0.833 & 0.008 & 0.830, 0.827, 0.842 & 0.839 [0.76, 0.91] & 123, 48, 62 & 0.0044 \\
    GO $+$ pathway multi-hots $+$ ProtT5 backbone $+$ FBA ratios & 0.849 & 0.005 & 0.852, 0.852, 0.844 & 0.851 [0.78, 0.92] & 131, 48, 126 & 0.0019 \\
    GO multi-hot $+$ pathway multi-hot $+$ ProtT5, the four ratios only (round 47 control) & 0.853 & 0.017 & 0.850, 0.871, 0.838 & 0.853 [0.78, 0.92] & 179, 69, 144 & 0.0027 \\
    GO multi-hot $+$ ProtT5 (no pathway block), the log(1 $+$ |v|) vector in the affine & 0.844 & 0.009 & 0.851, 0.834, 0.847 & 0.843 [0.76, 0.92] & 59, 51, 94 & 0.0031 \\
    GO multi-hot $+$ pathway multi-hot $+$ ProtT5, the log(1 $+$ |v|) vector in the affine, no layer & 0.869 & 0.010 & 0.867, 0.862, 0.880 & 0.873 [0.81, 0.93] & 137, 68, 185 & 0.0026 \\
    \hline
  \end{tabular}
\end{table}


\tcfigfit[1.0]{figures/w63_readout_2026-09-18-18-21-21.pdf}{\textbf{Round 47, the GO
  multi-hot beside the layer-free solve in one transformer.} \textbf{a}, Validation
  AUROC per epoch, one thin line per seed and the seed mean per arm. \textbf{b}, The
  median absolute epoch-to-epoch step per cell against the validation size.
  \textbf{c}, Held-out AUROC per cell (points), the seed mean and the gene bootstrap
  of the pooled per-gene logits (square with interval), beside FBA on four media,
  Flux Cone Learning and the released FBA predictions on the same 195 genes.
  \textbf{d}, Gradient norm per module per epoch, seed mean. Written by
  \texttt{experiments/028-gene-essentiality/scripts/wave12\_readout.py --round w63}.}{fig:w63}

GO enters the transformer on the logit at no cost and adds nothing there either
(round 50, jobs 2440, 2444 and 2446; Table~\ref{tab:w66}, Figure~\ref{fig:w66}). The
affine reads any store written in the solve store's shape, so the propagated GO
multi-hot was written as one (\texttt{go\_affine\_store.py}: 2,362 columns over the
1,121 labeled genes, 1,120 with an annotation; the trainer's $\log(1 + |v|)$ maps a
term to 0 or 0.693), alone and appended to the pFBA vector. GO alone as affine
columns beside pathway and ProtT5 on the node reads 0.886 on six seeds (sd 0.012;
pooled 0.892 [0.83, 0.94]) against 0.852 with GO on the node; GO with the
$\log(1 + |v|)$ vector as affine columns reads 0.895, 0.901, 0.906, 0.900, 0.893 and
0.911 (mean 0.901, sd 0.007; pooled 0.903 [0.85, 0.95]) against 0.866 with GO on the
node beside the same vector and 0.908 and 0.912 with the vector alone; the same on
ProtT5 alone 0.876 (pooled 0.880); and the GO-on-node cell run to 400 epochs 0.860
(pooled 0.864), no better than at 200. So the logit path carries the annotation
where the node path pays 0.04 for it, and the one model on the logit reads the
solve-alone number, not the two-family average (0.948). What the GO columns do in
the affine is visible in the training loss: with 2,362 binary columns over 629
training genes the cell's training loss crosses 0.6 at epoch 14 to 21 and reaches
0.01 to 0.04 by epoch 199, where the vector-alone cell crosses at epochs 33 to 58 and
ends near 0.1, and the selected epochs run 79 to 197 (251 to 379 at 400 epochs, which
reads 0.895); the affine memorizes the training genes through the GO columns and the
held-out ranking drifts on a fit plateau. Under the fold protocol (round 58, job
2453) the form reads 0.911, 0.907, 0.896, 0.886 and 0.884 on the folds and 0.918
[0.86, 0.97] as the ensemble on one seed, the solve-alone cell's 0.919 and 0.920
(Table~\ref{tab:ensembles}), so the one model on the logit is one number with the
solve alone on both protocols and 0.035 under the two-family average. A regularized
read of the GO columns is what remains of the one-model question on this backbone.

%% notes-tex/028-gene-essentiality/tables/round_w66.tex
%% GENERATED by experiments/028-gene-essentiality/scripts/notes_tex_tables.py
%% SOURCE: results/w66_readout.json
\begin{table}[H]
  \centering
  \scriptsize
  \caption[Round 50, GO on the logit instead of the node]{Round 50, GO on the logit instead of the node. The round-33 cell and the round-47 GO-on-node cells as the comparators, against the GO multi-hot (2,362 propagated terms) entering as affine columns through the solve store path, encoded as $\log(1 + |v|)$ (0 or 0.693 per term): GO alone beside the four ratios with pathway and ProtT5 on the node, GO appended to the pFBA vector (the one-model GO + layer-free form on the logit), the same on ProtT5 alone, and the round-47 GO-on-node cell run to 400 epochs. Batch 32, class weight; seeds 2700 to 2705 for the round-33 cell, 2700 to 2702 for the rest. Bold marks the highest value in a column. $\uparrow$ ($\downarrow$) marks an arm whose held-out AUROC is above (below) the first arm's on every one of the same three seeds.}
  \label{tab:w66}
  \begin{tabular}{@{}p{52mm}rrlllr@{}}
    \hline
    arm & mean & sd & seeds & pooled [95\% CI] & selected epochs & val step \\
    \hline
    the vector as log(1 $+$ |v|), default initialization & \textbf{0.908} & 0.008 & 0.896 to 0.919 (6) & \textbf{0.912} [0.87, 0.95] & 33 to 58 & 0.0043 \\
    GO multi-hot $+$ pathway multi-hot $+$ ProtT5, the four ratios only (round 47 control) & 0.852 & 0.011 & 0.838 to 0.871 (6) & 0.854 [0.78, 0.92] & 42 to 179 & 0.0028 \\
    GO multi-hot $+$ pathway multi-hot $+$ ProtT5, the log(1 $+$ |v|) vector in the affine, no layer & 0.866 & 0.009 & 0.857 to 0.880 (6) & 0.872 [0.81, 0.93] & 68 to 185 & 0.0026 \\
    the round-47 GO $+$ pathway $+$ ProtT5 layer-free cell at 400 epochs & 0.860 & 0.004 & 0.864, 0.862, 0.855 & 0.864 [0.80, 0.93] & 111, 255, 265 & 0.0015 \\
    pathway $+$ ProtT5 on the node; the GO multi-hot as affine columns (no solve) & 0.886 & 0.012 & 0.870 to 0.897 (6) & 0.892 [0.83, 0.94] & 83 to 194 & 0.0017 \\
    pathway $+$ ProtT5 on the node; the log(1 $+$ |v|) vector and the GO multi-hot as affine columns & 0.901 & 0.007 & 0.893 to 0.911 (6) & 0.903 [0.85, 0.95] & 79 to 197 & 0.0021 \\
    the same cell at 400 epochs & 0.895 & 0.006 & 0.892, 0.902, 0.890 & 0.895 [0.84, 0.94] & 251, 286, 379 & 0.0024 \\
    ProtT5 alone on the node; the log(1 $+$ |v|) vector and the GO multi-hot as affine columns & 0.876 & 0.003 & 0.874, 0.880, 0.875 & 0.880 [0.81, 0.94] & 164, 165, 176 & 0.0015 \\
    \hline
  \end{tabular}
\end{table}


\tcfigfit[1.0]{figures/w66_readout_2026-09-18-20-08-28.pdf}{\textbf{Round 50, GO on
  the logit instead of the node.} \textbf{a}, Validation AUROC per epoch, one thin
  line per seed and the seed mean per arm. \textbf{b}, The median absolute
  epoch-to-epoch step per cell against the validation size. \textbf{c}, Held-out AUROC
  per cell (points), the seed mean and the gene bootstrap of the pooled per-gene logits
  (square with interval), beside FBA on four media, Flux Cone Learning and the released
  FBA predictions on the same 195 genes. \textbf{d}, Gradient norm per module per
  epoch, seed mean. Written by
  \texttt{experiments/028-gene-essentiality/scripts/wave12\_readout.py --round w66}.}{fig:w66}

Weight decay is not the regularizer (round 59, job 2454; Table~\ref{tab:w75},
Figure~\ref{fig:w75}). The GO-and-vector affine cell with AdamW decay 0.1, a
thousand times the base $10^{-4}$, still crosses training loss 0.6 at epochs 15 to
16, ends at 0.03 to 0.05 and selects at epochs 173 to 182, reading 0.890, 0.898 and
0.891 (mean 0.893; pooled 0.892 [0.83, 0.94]) against 0.901 at the base; at decay 1.0
the cell selects at epochs 8 to 11 and reads 0.849, 0.849 and 0.845 (mean 0.848),
the decay pulling the affine back faster than the training genes can be fit. What the
columns need is fewer of them, or a read that cannot memorize 629 genes through
2,362 indicators; the reduced form is round 60.

%% notes-tex/028-gene-essentiality/tables/round_w75.tex
%% GENERATED by experiments/028-gene-essentiality/scripts/notes_tex_tables.py
%% SOURCE: results/w75_readout.json
\begin{table}[H]
  \centering
  \scriptsize
  \setlength{\tabcolsep}{2.5pt}
  \caption[Round 59, weight decay on the GO affine columns]{Round 59, weight decay on the GO affine columns. The round-33 cell (the vector alone in the affine) and the round-50 cell (GO and the vector as affine columns, AdamW weight decay $10^{-4}$) as the comparators, against the round-50 cell with weight decay 0.1 and 1.0. Batch 32, class weight, 200 epochs; seeds 2700 to 2705 for the comparators, 2700 to 2702 for the rest. Bold marks the highest value in a column. $\uparrow$ ($\downarrow$) marks an arm whose held-out AUROC is above (below) the first arm's on every one of the same three seeds.}
  \label{tab:w75}
  \begin{tabular}{@{}>{\raggedright\hangindent=1em\hangafter=1\arraybackslash}p{83.4mm}rrlllr@{}}
    \hline
    arm & mean & sd & seeds & pooled [95\% CI] & selected epochs & val step \\
    \hline
    the vector as log1p, default initialization & \textbf{0.908} & 0.008 & 0.896 to 0.919 (6) & \textbf{0.912} [0.87, 0.95] & 33 to 58 & 0.0043 \\
    pathway $+$ ProtT5 on the node; the log1p vector and the GO multi-hot as affine columns & 0.901 & 0.007 & 0.893 to 0.911 (6) & 0.903 [0.85, 0.95] & 79 to 197 & 0.0021 \\
    the same cell with AdamW weight decay 0.1 & 0.893 & 0.004 & 0.890, 0.898, 0.891 & 0.892 [0.83, 0.94] & 182, 173, 173 & 0.0021 \\
    the same cell with AdamW weight decay 1.0 & 0.848 & 0.003 & 0.849, 0.849, 0.845 & 0.848 [0.75, 0.93] & 10, 11, 8 & 0.0006 \\
    \hline
  \end{tabular}
\end{table}


\tcfigfit[1.0]{figures/w75_readout_2026-09-18-20-41-22.pdf}{\textbf{Round 59, weight
  decay on the GO affine columns.} \textbf{a}, Validation AUROC per epoch, one thin
  line per seed and the seed mean per arm. \textbf{b}, The median absolute
  epoch-to-epoch step per cell against the validation size. \textbf{c}, Held-out AUROC
  per cell (points), the seed mean and the gene bootstrap of the pooled per-gene logits
  (square with interval), beside FBA on four media, Flux Cone Learning and the released
  FBA predictions on the same 195 genes. \textbf{d}, Gradient norm per module per
  epoch, seed mean. Written by
  \texttt{experiments/028-gene-essentiality/scripts/wave12\_readout.py --round w75}.}{fig:w75}

Fewer columns stop the memorizing and read the solve-alone number (round 60, jobs
2455 and 2458; Table~\ref{tab:w76}, Figure~\ref{fig:w76}). The GO multi-hot
projected on its first 64 principal components fitted over the 1,121 labeled genes
(80 percent of the variance), split into positive and negative parts so that the
$\log(1 + |v|)$ encoding keeps the sign (128 columns), reads 0.886 on six seeds
alone (sd 0.022; pooled 0.902 [0.85, 0.95]) and 0.905, 0.911, 0.919, 0.900, 0.900 and
0.902 with the vector (mean 0.906, sd 0.008; pooled 0.913 [0.86, 0.96]), the
vector-alone cell's 0.908 and 0.912, the cells crossing training loss 0.6 at epochs
40 to 66 and selecting at 33 to 74 as the vector-alone cell does; under the fold
protocol (round 62, jobs 2459 to 2461) the same cell reads 0.914 [0.86, 0.96] on
one seed and 0.917 [0.87, 0.96] on three (folds 0.917, 0.878, 0.843, 0.901 and
0.912 as seed-mean logits), the solve-alone 0.919 and 0.920
(Table~\ref{tab:ensembles}). Sixteen components (round 63, jobs 2462 and 2467; 32
columns, 52 percent of the variance) read the same: 0.874 alone (sd 0.029; pooled
0.893) and 0.908 with the vector (sd 0.016; pooled 0.915 [0.86, 0.96]) on six seeds.
The one model on this backbone is closed on the split file and the folds: GO on the
node costs 0.04, GO on the logit costs nothing and adds nothing, in full it memorizes
and reduced to 64 or 16 components it reads the solve alone, so the annotation's 0.04
over the two families' average is not the transformer's to take through either path,
and the two families stay two models.

%% notes-tex/028-gene-essentiality/tables/round_w76.tex
%% GENERATED by experiments/028-gene-essentiality/scripts/notes_tex_tables.py
%% SOURCE: results/w76_readout.json
\begin{table}[H]
  \centering
  \scriptsize
  \caption[Round 60, the GO multi-hot reduced before the affine]{Round 60, the GO multi-hot reduced before the affine. The round-33 cell and the round-50 cells (the full GO multi-hot as affine columns, alone and with the vector) as the comparators, against the GO multi-hot projected on its first 64 principal components fitted over the 1,121 labeled genes and split into positive and negative parts (128 columns, so the $\log(1 + |v|)$ encoding keeps the sign), alone and with the vector. Batch 32, class weight, 200 epochs; seeds 2700 to 2705 for the comparators, 2700 to 2702 for the rest. Bold marks the highest value in a column. $\uparrow$ ($\downarrow$) marks an arm whose held-out AUROC is above (below) the first arm's on every one of the same three seeds.}
  \label{tab:w76}
  \begin{tabular}{@{}p{52mm}rrlllr@{}}
    \hline
    arm & mean & sd & seeds & pooled [95\% CI] & selected epochs & val step \\
    \hline
    the vector as log(1 $+$ |v|), default initialization & 0.908 & 0.008 & 0.896 to 0.919 (6) & 0.912 [0.87, 0.95] & 33 to 58 & 0.0043 \\
    pathway $+$ ProtT5 on the node; the GO multi-hot as affine columns (no solve) & 0.886 & 0.012 & 0.870 to 0.897 (6) & 0.892 [0.83, 0.94] & 83 to 194 & 0.0017 \\
    pathway $+$ ProtT5 on the node; the log(1 $+$ |v|) vector and the GO multi-hot as affine columns & 0.901 & 0.007 & 0.893 to 0.911 (6) & 0.903 [0.85, 0.95] & 79 to 197 & 0.0021 \\
    pathway $+$ ProtT5 on the node; GO reduced to 64 principal components (positive and negative parts) as affine columns & 0.896 & 0.003 & 0.895, 0.895, 0.899 & 0.904 [0.85, 0.95] & 33, 41, 74 & 0.0034 \\
    the same with the log(1 $+$ |v|) vector as affine columns too & \textbf{0.912} & 0.007 & 0.905, 0.911, 0.919 & \textbf{0.917} [0.87, 0.96] & 69, 53, 52 & 0.0024 \\
    \hline
  \end{tabular}
\end{table}


\tcfigfit[1.0]{figures/w76_readout_2026-09-18-22-17-30.pdf}{\textbf{Rounds 60 and 63,
  the GO multi-hot reduced before the affine.} \textbf{a}, Validation AUROC per epoch, one
  thin line per seed and the seed mean per arm. \textbf{b}, The median absolute
  epoch-to-epoch step per cell against the validation size. \textbf{c}, Held-out AUROC
  per cell (points), the seed mean and the gene bootstrap of the pooled per-gene logits
  (square with interval), beside FBA on four media, Flux Cone Learning and the released
  FBA predictions on the same 195 genes. \textbf{d}, Gradient norm per module per
  epoch, seed mean. Written by
  \texttt{experiments/028-gene-essentiality/scripts/wave12\_readout.py --round w76}.}{fig:w76}

\subsection{Round 30: the genome-wide labels as a larger held-out set\secstatus{todo}}
\label{sec:round30}

Every number above sits on 195 genes, and the open list's first question was whether
the graph convolution's margin survives a larger held-out set. The genome-wide labels
of round 1 give one: 4,674 genes the released file never labeled, 974 essential by SGD
annotation and 3,700 nonessential by a deletion-collection fitness record. Round 30
trains on the released labeled genes outside the released held-out set (the split
\texttt{ess\_genome\_heldout\_val0.2\_s0}: 629 train, 158 validation, the 195 released
held-out genes dropped so the read is not on them) and reads once on the 4,674
(Table~\ref{tab:genome-heldout}). Two properties of the set bound what it can say.
None of the 4,674 carries a yeast-GEM reaction, since the released file had labeled the
GEM's genes, so the read is a transfer from metabolic to non-metabolic genes and FBA's
knockout ratio is 1.0 on every one of them, an AUROC of exactly 0.5. And label origin
coincides with class: every essential gene is an SGD annotation and every nonessential
one a fitness record, so a model that can tell the two sources apart scores on this set
without reading essentiality.

%% notes-tex/028-gene-essentiality/tables/genome_heldout.tex
%% GENERATED by experiments/028-gene-essentiality/scripts/notes_tex_tables.py
%% SOURCE: results/genome_heldout_readout.json
\begin{table}[H]
  \centering
  \scriptsize
  \caption[Round 30, the genome-wide held-out genes]{Round 30, the genome-wide labels as a larger held-out set. Every model is trained on the released labeled genes outside the released held-out set (629 train, 158 val, the split \texttt{ess\_genome\_heldout\_val0.2\_s0}) and read once on the 4,674 genes the released file never labeled, 974 of them essential by SGD annotation and the rest nonessential by a fitness record, so label origin coincides with class on this set. None of these genes carries a yeast-GEM reaction, so FBA's knockout ratio is 1.0 on every one and the solve has no call. Cells are seeds 2700 to 2702; the pooled column scores the seed-mean logit with a 2,000-draw gene bootstrap. Bold marks the highest value in a column.}
  \label{tab:genome-heldout}
  \begin{tabular}{@{}p{58mm}rlrrlp{22mm}@{}}
    \hline
    arm & cells & per cell & mean & sd & pooled [95\% CI] & selected epochs \\
    \hline
    pathway multi-hot + ProtT5 + FBA affine (round 22 recipe) & 6 & 0.571, 0.571, 0.579, 0.569, 0.571, 0.573 & 0.572 & 0.003 & 0.572 [0.55, 0.59] & -- \\
    cached solve + flux head, pathway + ProtT5 (round 27 form) & 6 & 0.543, 0.480, 0.531, 0.512, 0.579, 0.543 & 0.532 & 0.033 & 0.551 [0.53, 0.57] & -- \\
    layer-free log1p vector + affine, pathway + ProtT5 (round 33 form) & 6 & 0.542, 0.561, 0.569, 0.557, 0.568, 0.599 & 0.566 & 0.019 & 0.570 [0.55, 0.59] & -- \\
    MLP on ProtT5 alone & 3 & 0.539, 0.531, 0.532 & 0.534 & 0.003 & 0.534 [0.51, 0.55] & 62, 58, 56 \\
    MLP on ProtT5 + FBA ratios & 3 & 0.494, 0.476, 0.488 & 0.486 & 0.008 & 0.481 [0.46, 0.50] & 71, 197, 76 \\
    MLP on the three inputs & 3 & 0.511, 0.508, 0.507 & 0.509 & 0.002 & 0.509 [0.49, 0.53] & 55, 58, 55 \\
    graph convolution, physical graph, all-ones input & 3 & 0.684, 0.683, 0.692 & 0.686 & 0.004 & 0.686 [0.67, 0.70] & 91, 100, 77 \\
    graph convolution, physical graph, degree and clustering & 3 & 0.757, 0.752, 0.749 & 0.753 & 0.003 & 0.759 [0.74, 0.78] & 28, 31, 3 \\
    graph convolution, physical graph, ProtT5 & 3 & 0.545, 0.542, 0.544 & 0.544 & 0.001 & 0.544 [0.52, 0.57] & 192, 197, 188 \\
    graph convolution, union of nine, all-ones input & 3 & 0.643, 0.643, 0.643 & 0.643 & 0.000 & 0.643 [0.62, 0.66] & 20, 5, 3 \\
    graph convolution, union of nine, degree and clustering & 3 & 0.660, 0.667, 0.660 & 0.663 & 0.003 & 0.663 [0.64, 0.68] & 197, 196, 197 \\
    graph convolution, union of nine, ProtT5 & 3 & 0.755, 0.753, 0.756 & 0.755 & 0.001 & 0.755 [0.74, 0.77] & 197, 197, 197 \\
    graph attention, physical graph, degree and clustering & 3 & 0.709, 0.692, 0.720 & 0.707 & 0.011 & 0.708 [0.69, 0.73] & 179, 190, 119 \\
    graph attention, union of nine, degree and clustering & 3 & 0.715, 0.704, 0.719 & 0.713 & 0.006 & 0.717 [0.70, 0.74] & 197, 197, 197 \\
    graph convolution on ProtT5 alone, 400 epochs & 3 & 0.746, 0.753, 0.749 & 0.749 & 0.003 & 0.749 [0.73, 0.77] & 373, 380, 393 \\
    graph convolution on ProtT5 + FBA ratios & 3 & 0.759, 0.758, 0.758 & \textbf{0.759} & 0.000 & 0.759 [0.74, 0.78] & 196, 197, 197 \\
    graph convolution on the three inputs, 400 epochs & 3 & 0.742, 0.744, 0.744 & 0.744 & 0.001 & 0.744 [0.72, 0.76] & 380, 396, 380 \\
    graph convolution on GO alone, 400 epochs & 3 & 0.716, 0.697, 0.694 & 0.702 & 0.010 & 0.703 [0.68, 0.72] & 397, 395, 379 \\
    graph convolution on GO + FBA ratios, 400 epochs & 3 & 0.699, 0.693, 0.707 & 0.700 & 0.006 & 0.700 [0.68, 0.72] & 397, 392, 397 \\
    graph convolution on ProtT5 + GO + FBA ratios, 400 epochs & 3 & 0.700, 0.697, 0.707 & 0.701 & 0.004 & 0.702 [0.68, 0.72] & 397, 368, 391 \\
    FBA knockout ratio on YPD (0 of 4674 genes carry a reaction) & 1 & 0.500 & 0.500 & 0.000 & 0.500 [0.50, 0.50] & -- \\
    \hline
  \end{tabular}
\end{table}


%% notes-tex/028-gene-essentiality/tables/degree_by_class.tex
%% GENERATED by experiments/028-gene-essentiality/scripts/notes_tex_tables.py
%% SOURCE: results/degree_by_class.json
\begin{table}[H]
  \centering
  \scriptsize
  \caption[Degree as the score]{The gene's degree in each graph as the essentiality score, nothing trained, on the 4,674 genome-wide held-out genes of round 30, the 195 released held-out genes, and the 787 released genes the models train and validate on. AUROC of degree with a 2,000-draw gene bootstrap, and the median degree of the essential against the nonessential genes in that set. Bold marks the highest AUROC in a column. Written by \texttt{experiments/028-gene-essentiality/scripts/degree\_by\_class.py}.}
  \label{tab:degree-by-class}
  \begin{tabular}{@{}lllllll@{}}
    \hline
    & \multicolumn{2}{l}{genome held-out (4,674)} & \multicolumn{2}{l}{released held-out (195)} & \multicolumn{2}{l}{released train + val (787)} \\
    graph & AUROC & median ess / non & AUROC & median ess / non & AUROC & median ess / non \\
    \hline
    physical & \textbf{0.740} [0.72, 0.76] & 50 / 16 & 0.594 [0.47, 0.71] & 27 / 19 & 0.591 [0.53, 0.65] & 21 / 17 \\
    regulatory & 0.478 [0.46, 0.50] & 4 / 5 & 0.513 [0.41, 0.62] & 6 / 7 & 0.474 [0.42, 0.53] & 6 / 6 \\
    tflink & 0.562 [0.55, 0.58] & 30 / 27 & 0.424 [0.31, 0.53] & 37 / 46 & 0.424 [0.37, 0.48] & 37 / 42 \\
    STRING neighborhood & 0.514 [0.50, 0.53] & 0 / 0 & \textbf{0.618} [0.50, 0.73] & 175 / 67 & \textbf{0.629} [0.57, 0.68] & 142 / 68 \\
    STRING fusion & 0.625 [0.61, 0.64] & 1 / 0 & 0.605 [0.49, 0.71] & 7 / 4 & 0.555 [0.50, 0.61] & 6 / 5 \\
    STRING cooccurence & 0.497 [0.48, 0.51] & 0 / 0 & 0.356 [0.27, 0.45] & 0 / 1 & 0.368 [0.32, 0.42] & 0 / 1 \\
    STRING coexpression & 0.681 [0.66, 0.70] & 394 / 182 & 0.617 [0.49, 0.73] & 477 / 292 & 0.585 [0.52, 0.64] & 385 / 300 \\
    STRING experimental & 0.638 [0.62, 0.66] & 263 / 121 & 0.594 [0.47, 0.70] & 203 / 158 & 0.517 [0.46, 0.57] & 139 / 149 \\
    STRING database & 0.714 [0.70, 0.73] & 22 / 1 & 0.577 [0.48, 0.68] & 41 / 29 & 0.485 [0.43, 0.54] & 24 / 29 \\
    union of nine & 0.655 [0.64, 0.67] & 534 / 297 & 0.611 [0.49, 0.73] & 660 / 412 & 0.572 [0.51, 0.63] & 540 / 449 \\
    \hline
  \end{tabular}
\end{table}


The inputs do not transfer and the graph does. The MLP on the transformer's three
inputs reads 0.509 on the 4,674 genes (0.486 on ProtT5 with the ratios), chance, from
0.870 on the released 195; the same inputs through two layers of graph convolution read
0.744, ProtT5 alone through the graph 0.749 and with the ratios 0.759, and GO alone
through the graph 0.702, under the three-input convolution here where on the 195 it
was 0.02 above it. The graph convolution on GO does not keep its margin off the
metabolic genes; the graph itself carries what transfers, and the ratios add 0.01 to a
convolution on ProtT5 when every ratio in the test set is 1.0, which measures the
training set and not the test set. What the graph carries is the gene's neighborhood
in the physical graph, measured directly with the ladder's structure-only inputs
(job 2401): degree and clustering alone, through two layers of convolution over the
physical graph, read 0.753 [0.74, 0.78], level with ProtT5 through the union graph
(0.755) and with the best number on the set, while the all-ones input reads 0.686 on
the physical graph and 0.643 on the union, ProtT5 without a graph 0.534, and ProtT5
through the physical graph at 200 epochs 0.544. On these genes the number is the
degree of the gene in the protein interaction graph, which a convolution reads from
any input that lets it count neighbors and a model without the graph cannot; whether
that is essentiality or the deletion collection's own coverage the label origin does
not let this set decide. Degree needs no model to say so
(Table~\ref{tab:degree-by-class}): the gene's degree in the physical graph, ranked as
it is, reads 0.740 [0.72, 0.76] on the 4,674 genes, a median of 50 interactions for an
essential gene against 16 for a nonessential one, and 0.594 [0.47, 0.71] on the
released 195 and 0.591 on the 787 released training genes, where the medians are 27
against 19 and 21 against 17. The STRING database channel's degree reads 0.714 on the
genome-wide genes and 0.577 on the released ones, the union's 0.655 and 0.611, and no
other graph's degree passes 0.69 anywhere. So the physical graph's degree is a
genome-wide essentiality signal that the metabolic genes barely vary in, which is why
no model in this document read it on the 195 and every graph model read it on the
4,674. The transformer does not transfer either: the
round-22 recipe, the pathway multi-hot with ProtT5 and the FBA affine, trained on the
same 629 genes reads 0.572 [0.55, 0.59] on the 4,674 (0.569 to 0.579 over six seeds,
jobs 2396 and 2544), 0.07 above chance and 0.18 under the graph convolution on ProtT5
alone,
from 0.890 on the released 195. Its inputs are the MLP's, and in this recipe the gene
graphs enter training only through the graph regularizer, not as a message-passing
input, so the three models that read 0.87 to 0.89 on the metabolic genes without the
graph as an input all fall to chance or near it off them, and the one that convolves
over the graph does not. The scalable module form on the same split (jobs 2397, 2545 and 2546, six
seeds) reads 0.551 [0.53, 0.57] pooled, and is the one arm here whose seeds disagree:
0.480 to 0.579, a spread of 0.033 against the affine recipe's 0.003, with one seed
below chance. It sits under the affine recipe: its solved flux is the wild-type vector on every one of these genes, a
constant the head cannot rank by, so what it adds over the backbone is nothing by
construction, and what it reads is the backbone with a head that fit the metabolic
genes' flux. The layer-free form (round 65, jobs 2464 and 2466, six seeds) reads
0.566 (sd 0.019; pooled 0.570 [0.55, 0.59]) on the same split, between the two, for
the same reason: the solve store covers the 1,121 labeled genes only, so the affine's solve columns read one all-zero row on every held-out gene (\texttt{n\_genes\_with\_solve} 1,121 in the result files), a constant that shifts every logit alike.

The other direction says the larger label set is not the missing training data for
the metabolic genes (Tables~\ref{tab:graph-go400-genome}
and~\ref{tab:graph-extra-genome}). Trained on the 4,368 genome-labeled genes and read
on the released 195, the graph convolution on GO reads 0.913 on every seed (0.912,
0.913, 0.913), the number it reads from 629 genes with the seed spread closed, and the
cheap models on the transformer's three inputs lose to their 629-gene runs: the MLP
0.837 from 0.870, the 400-epoch convolution 0.843 from 0.888. Seven times the labels,
of a different origin for six sevenths of them, buy the GO convolution nothing and
cost the three-input models 0.03 to 0.05, which is what round 1 measured on the
transformer.

%% notes-tex/028-gene-essentiality/tables/graph_baselines_go400_genome.tex
%% GENERATED by experiments/028-gene-essentiality/scripts/notes_tex_tables.py
%% SOURCE: results/graph_baselines_go400_genome.json
\begin{table}[H]
  \centering
  \footnotesize
  \caption[Graph convolution on GO trained on the genome-wide labels]{The three rows of Table~\ref{tab:graph-go400} trained on the genome-wide labels (split \texttt{ess\_genome\_val0.2\_s0}: 4,368 train and 1,093 validation genes, the released labels where the file has them and SGD or a fitness record elsewhere) and read on the same 195 released held-out genes, seeds 2700 to 2702. Bold marks the highest value in a column.}
  \label{tab:graph-go400-genome}
  \begin{tabular}{@{}lllrrrrll@{}}
    \hline
    model & graph & inputs & dim & params & mean & sd & seeds & selected epochs \\
    \hline
    GCN & union of nine & GO (400 epochs) & 2,362 & 155,457 & \textbf{0.913} & 0.000 & 0.912, 0.913, 0.913 & 397, 397, 397 \\
    GCN & union of nine & GO $+$ FBA ratios (400 epochs) & 2,366 & 155,713 & 0.911 & 0.002 & 0.912, 0.914, 0.908 & 397, 397, 395 \\
    GCN & union of nine & ProtT5 $+$ GO $+$ FBA ratios (400 epochs) & 3,390 & 221,249 & 0.907 & 0.002 & 0.906, 0.910, 0.906 & 396, 397, 397 \\
    \hline
  \end{tabular}
\end{table}


%% notes-tex/028-gene-essentiality/tables/graph_baselines_extra_genome.tex
%% GENERATED by experiments/028-gene-essentiality/scripts/notes_tex_tables.py
%% SOURCE: results/graph_baselines_extra_genome.json
\begin{table}[H]
  \centering
  \scriptsize
  \caption[Cheap models on the three inputs trained on the genome-wide labels]{The MLP and graph convolution on the transformer's inputs (Table~\ref{tab:graph-baselines}) trained on the genome-wide labels (split \texttt{ess\_genome\_val0.2\_s0}) and read on the same 195 released held-out genes, seeds 2700 to 2702. Bold marks the highest value in a column.}
  \label{tab:graph-extra-genome}
  \begin{tabular}{@{}llp{30mm}rrrrll@{}}
    \hline
    model & graph & inputs & dim & params & mean & sd & seeds & selected epochs \\
    \hline
    MLP & none & ProtT5 $+$ chrom-pathway & 1,221 & 82,433 & 0.808 & 0.004 & 0.812, 0.803, 0.810 & 101, 71, 71 \\
    GCN & union of nine & ProtT5 $+$ chrom-pathway & 1,221 & 82,433 & 0.824 & 0.001 & 0.825, 0.825, 0.823 & 197, 197, 197 \\
    MLP & none & ProtT5 $+$ FBA ratios & 1,028 & 70,081 & 0.770 & 0.000 & 0.769, 0.770, 0.770 & 79, 80, 79 \\
    GCN & union of nine & ProtT5 $+$ FBA ratios & 1,028 & 70,081 & 0.818 & 0.002 & 0.819, 0.820, 0.815 & 197, 196, 197 \\
    MLP & none & ProtT5 $+$ chrom-pathway $+$ FBA ratios & 1,225 & 82,689 & 0.837 & 0.002 & 0.837, 0.840, 0.834 & 68, 80, 71 \\
    GCN & union of nine & ProtT5 $+$ chrom-pathway $+$ FBA ratios & 1,225 & 82,689 & 0.824 & 0.003 & 0.822, 0.822, 0.827 & 197, 197, 197 \\
    GCN & union of nine & ProtT5 (400 epochs) & 1,024 & 69,825 & 0.831 & 0.002 & 0.828, 0.831, 0.833 & 397, 397, 396 \\
    GCN & union of nine & ProtT5 $+$ chrom-pathway $+$ FBA ratios (400 epochs) & 1,225 & 82,689 & \textbf{0.843} & 0.002 & 0.841, 0.843, 0.846 & 395, 395, 395 \\
    \hline
  \end{tabular}
\end{table}

